These notes are an introduction to probabilistic thinking: the mathematical art of reasoning under uncertainty. Probability is not presented here as a collection of ready-made formulas, but as a language for describing experiments, incomplete information, variability, randomness, and data.

The purpose of the course is to teach how to build probabilistic models carefully. Before one computes a probability, expectation, variance, confidence interval, or approximation, one should understand what is random, what is being observed, what assumptions are being made, and what mathematical object represents the situation. In this sense, probability is not only a computational tool, but also a discipline of precise modelling.

The course is intended for readers who want to understand why probabilistic methods work, not only how to apply them. The emphasis is therefore placed on interpretation, model construction, and the connection between formal definitions and concrete situations. Calculations appear throughout the notes, but they are treated as consequences of a well-defined model rather than as isolated procedures.

\section{How to Use These Notes}

These notes are meant to be read before problem classes. Their role is not only to present the material, but also to show how mathematical argumentation is built: how one starts from a situation, identifies the relevant objects, formulates assumptions, chooses notation, develops a line of reasoning, and only then performs computations. A problem class is therefore not the place where the basic language of the topic is introduced for the first time; it is the place where that language is used, tested, discussed, and refined through problems.

Students are expected to read the relevant chapter in advance and to treat the text as a model of reasoning. The important question is not only ``What is the answer?'' but also ``How is the solution argued?'' While reading, students should pay attention to the order of the argument, the meaning of each definition, the reason for each assumption, and the way examples are connected to general concepts. The aim is to learn how to construct a probabilistic explanation, not merely how to recognize a formula that seems to fit a task.

Problem lists will accompany the chapters and will be used during classes. They are meant to be solved after the corresponding text has been read. The chapter gives the language, concepts, and patterns of argumentation needed to approach the problems responsibly. Without this preparation, a solution may look complete on paper while still being mathematically empty: it may contain formulas, calculations, or generated text, but no clear understanding of what model is being used or why the reasoning is valid.

Artificial intelligence tools may be useful for practice, checking calculations, generating additional examples, or asking for alternative explanations. Students are encouraged to use such tools responsibly when they support learning. However, AI-generated text cannot replace the student's own reading and understanding. A student must be able to read, question, correct, and explain any solution they present. If one cannot explain the experiment, the elementary outcomes, the assumptions, the chosen model, and the steps of the argument, then the work has not yet been done.

These notes also help students learn what to ask for. To use problem lists, textbooks, or AI tools well, one must first know what matters: which object is being modelled, which assumptions are decisive, which definitions are involved, and what kind of justification is expected. The PDF is therefore an essential part of the course: it teaches the material, but it also teaches how to read mathematics, how to demand proper explanations, and how to recognize whether an answer contains real reasoning.

The expected preparation for classes is conceptual as well as computational. After reading a chapter, a student should be able to summarize its main ideas, explain the definitions in their own words, reconstruct the key examples, identify the assumptions in a probabilistic model, follow and reproduce the main arguments, and answer questions that test understanding of the text. The notes are deliberately written with full explanations and detailed argumentation so that this preparation is possible.

\section{Structure of the Course}

These lecture notes are organized as a gradual construction of probabilistic thinking. The aim is not to begin with formulas and then search for examples, but to develop the mathematical language needed to describe uncertainty in a precise way. We start from the problem of counting and representing possible outcomes, then move toward events, probability, random variables, probability laws, empirical data, sampling, limit theorems, and statistical inference.

The course begins with combinatorics because many probabilistic models depend on the ability to count possible outcomes correctly. At this stage probability itself is not yet needed. What is needed is the more elementary skill of recognizing what kind of object an outcome is: a sequence, a selection, an arrangement, a choice with or without repetition, or an object whose elements may or may not be distinguishable. This first step builds the counting language that later allows us to describe sample spaces and favorable outcomes with precision.

The second stage has a special conceptual role. It is devoted to discovering the formal language of probability theory. The purpose is to show how a mathematical theory can be built from careful thinking, analysis, and abstraction. We examine how statements about a random situation select sets of outcomes, how these sets become events, how logical operations correspond to set operations, and why a numerical theory of probability should be placed on such a structure. This part is not primarily computational; it is meant to teach how probabilistic language is formed.

After this conceptual preparation, the notes turn to sample spaces and elementary probability. Several canonical examples are developed from the beginning: one first understands the experiment, decides what should count as an elementary outcome, constructs the sample space, defines events, and only then assigns probabilities. Different representations of the same structure, such as lists, tables, trees, and symbolic notation, are used to show that a sample space is an abstract mathematical model rather than a collection of observed data.

The next step is conditional probability. Here the central idea is that information changes the space in which we reason. Conditional probability is introduced by asking what it means to restrict attention to the cases in which some event has already occurred. From this point of view, the formula for conditional probability becomes natural, and it leads to independence, dependence, partitions, the law of total probability, and Bayes' theorem. Bayes' theorem is treated as a central tool for reversing conditional reasoning and for understanding situations where intuition is easily misled.

The course then introduces random variables through the discrete case. A random variable appears as a function that extracts a numerical quantity from an elementary outcome: the number of successes, the number of errors, the number of arrivals, a payoff, or another measurable feature of the experiment. This leads to the support of a discrete random variable, its probability mass function, its cumulative distribution function, expectation, variance, moments, and interpretation. Classical discrete laws are presented as models arising from specific mechanisms, not as isolated formulas.

The continuous case is then developed as a parallel but conceptually more demanding part of the theory. For a continuous random variable, probability is no longer concentrated at individual points, so probabilities of intervals become the fundamental objects. This leads to density functions, cumulative distribution functions, integrals, continuous expectations, variances, moments, quantiles, and the interpretation of probability as area under a density. Important continuous models, including the uniform, exponential, normal, gamma, and beta distributions, are introduced with attention to their meaning and use.

Once theoretical probability laws have been introduced, the course moves to empirical distributions and descriptive statistics. The same ideas now appear in data: frequency tables, histograms, empirical cumulative distribution functions, means, medians, variances, standard deviations, quartiles, interquartile ranges, and outliers. This stage is meant to clarify the difference between an idealized probability law and a finite dataset, and also to show how empirical summaries mirror the theoretical quantities defined earlier.

The following stage adds another level of randomness by studying sampling and distributions of statistics. A sample mean, sample proportion, sample variance, median, or other statistic is not merely a number; it is itself a random variable because it depends on the random sample. This perspective prepares the ground for understanding sampling distributions and for seeing why repeated sampling creates a new layer of probabilistic structure.

The final stage of the present course outline is devoted to limit theorems and statistical inference. Once empirical summaries and sampling distributions are understood, we can ask what regularities appear as the sample size grows or as experiments are repeated many times. This leads to the law of large numbers, the central limit theorem, normal approximations, standard errors, confidence intervals, and the basic logic of statistical inference. The goal is to show how probability theory makes inference from data possible: it allows us to estimate unknown quantities, measure uncertainty, and justify conclusions mathematically.
